\section{Conclusion and Future Work}
\label{bab:5}

This research designed and evaluated a network-latency-aware eviction filtering architecture for the Kubernetes Descheduler. Our evaluation demonstrates that network-aware filtering can be effectively integrated via a dual-layer architecture, combining a global pre-eviction filter with strategy-level assessments using real-time telemetry or topology labels. This architecture preserves communication efficiency without compromising resource-balancing objectives, provided non-network-sensitive pods are available to sacrifice. Specifically, the system successfully capped latency spikes while preventing all cross-region pod spills, avoiding the significant latency penalties observed in upstream baselines. However, our evaluation also revealed a critical systemic limitation: isolated descheduler intent can be subverted by the network-blind default Kubernetes scheduler, which frequently misrouted evicted pods to degraded nodes based on raw compute availability. Ultimately, our findings prove that network protection and geometric resource balance can be achieved together when alternative eviction candidates exist. Future work includes establishing explicit coordination mechanisms with the primary kube-scheduler, weighting network costs using service-mesh traffic volumes, and integrating eviction logic directly with native Kubernetes Topology Spread Constraints to provide a more holistic approach to cluster network balancing.
